% !TEX root = tarasov-qasp-zh.tex

\chapter[你知道光束是怎样反射和折射的吗？]{你知道光束是怎样反射和折射的吗？\label{ch-31}}

\begin{dialogue}

\tchr 请陈述光的反射定律和折射定律。

\stda 反射定律是：入射角等于反射角。折射定律：入射角的正弦与折射角的正弦之比，等于该介质的折射率。
\begin{figure}[!h]%
\centering
\includegraphics[width=0.4\textwidth]{figs/ch-31/fig-128a}
\caption[反射可能在不同平面内发生。]{反射可能在不同平面内发生。\label{fig-128}}
\end{figure}
\tchr 你的这些表述很不准确。首先，你根本没有提到：入射光线和反射（或折射）光线，与在入射点处所作的、垂直于反射（或折射）界面的法线，位于同一平面内。如果不指明这一点，就可能以为反射是像\figr{fig-128} 所画的那样进行的。其次，你对折射定律的表述只涉及光线从空气射到某种介质界面上这一特殊情形。设一般情形下，光线从折射率为 $n_{1}$ 的介质射到折射率为 $n_{2}$ 的介质的界面上。入射角记作 $\alpha_{1}$，折射角记作 $\alpha_{2}$。这时折射定律可以写成
\begin{equation}
\frac{\sin \alpha_{1}}{\sin \alpha_{2}} = \frac{n_{2}}{n_{1}}
\label{eq-189}
\end{equation}
只要认为空气的 $n_{1}  = 1$，就得到你刚才的表述。

再来看下面的问题。\hlred{一枚硬币位于水中深度为 $H$ 处。我们沿竖直方向从上往下看它。我们会看到硬币位于多大的深度？}
\pnote{“三线共面”（入射线、法线、反射线在同一平面内）与“两角相等”合起来才是完整的反射定律，见人教版初中物理八年级上册第四章第 2 节《光的反射》；高中在选择性必修第一册第四章第 1 节《光的折射》给出折射定律 $n_{1}\sin\theta_{1}=n_{2}\sin\theta_{2}$，并定义折射率 $n=c/v$（真空光速与介质中光速之比）。式~\eqref{eq-189} 右边的 $n_{2}/n_{1}$ 即相对折射率；学生开头“正弦之比等于介质折射率”的说法只在光从真空（空气）射入介质时成立，这正是教师指出的第二点。}

\stda 我知道硬币看上去会浮起一些。更确切的答案，恐怕我给不出来。
\begin{figure}[h]%
\centering
\includegraphics[width=.6\textwidth]{figs/ch-31/fig-129a}
\caption[硬币看起来会位于何处？]{硬币看起来会位于何处？\label{fig-129}}
\end{figure}
\tchr 从硬币的中心作两条光线：$OA$ 和 $OB_{1}B$（\figr{fig-129}）。光线 $OA$ 是竖直的，不发生折射；光线 $OB_{1}B$ 则发生折射。设这两条发散的光线进入眼睛。眼睛将在发散光线 $OA$ 与 $B_{1}B$ 的交点处，即点 $O_{1}$ 处，看到硬币的像。由图显然可见，所求距离 $h$ 与深度 $H$ 之间有如下关系
\begin{equation*}%
h \tan \alpha_{1} = H \tan \alpha_{2}
\end{equation*}
由此得到
\begin{equation}%
h =H \, \frac{\tan \alpha_{2}}{\tan \alpha_{1}}
\label{eq-190}
\end{equation}
由于角 $\alpha_{1}$ 和 $\alpha_{2}$ 都很小，可以应用近似公式
 \begin{equation}
\tan \alpha \approx \sin \alpha \approx \alpha
\label{eq-191}
\end{equation}
（式中角度以弧度为单位，而不是以度为单位）。利用公式~\eqref{eq-191}，可以把方程~\eqref{eq-190} 改写成如下形式：
\begin{equation}
h \approx H \, \frac{\sin \alpha_{2}}{\sin \alpha_{1}} = \frac{H}{n}
\label{eq-192}
\end{equation}
水的 $n = 4/3$，于是
\begin{equation*}
h = \frac{3}{4}H
\end{equation*}
\pnote{$h=H/n$ 是小角度近似下的结果：只有当眼睛几乎竖直下看（$\alpha_{1}$、$\alpha_{2}$ 都很小）时，才能用 $\tan\alpha\approx\alpha$ 把式~\eqref{eq-190} 化简成式~\eqref{eq-192}。这就是“池水变浅”现象的定量形式——水深 2~m 看起来约 1.5~m（人教版初中物理八年级上册第四章第 4 节《光的折射》定性讨论过这一现象）。斜视时像既浮高又侧移，即下文所说“移远”。}

\stdb 如果我们不是竖直地、而是从侧面看这枚硬币，结果会怎样？

\tchr 这时硬币看上去不仅浮高了，而且移远了（见\figr{fig-129} 中的虚线）。显然，这种情形下的计算要复杂得多。再来看下面的问题。\hlred{一名身高为 $h$ 的潜水员站在深度为 $H$ 的湖底上。试求从他所站的点到那些他能经水面反射看到的湖底点的最小距离。}
\begin{figure}[h]%
\centering
\includegraphics[width=.6\textwidth]{figs/ch-31/fig-130a}
\caption[求到经水面反射可见的湖底点的最小距离。]{求从潜水员所站的点到他能经水面反射看到的湖底点的最小距离。\label{fig-130}}
\end{figure}
\stda 这类问题我知道怎么解。设所求距离为 $L$。从点 $A$ 到潜水员眼睛的光路如\figr{fig-130} 所示。点 $A$ 是潜水员能经湖面反射看到的最靠近他的点。比如，从更近的点 $B$ 出发的光线，经水面折射后掠过水面射出，不会返回潜水员眼中（见\figr{fig-130} 中的虚线）。角 $\alpha$ 是全反射（total internal reflection）的临界角（critical angle），由下式求得：
\begin{equation}%
\sin \alpha = \frac{1}{n}
\label{eq-193}
\end{equation}
由图可见
 \begin{align*}%
L & = h \tan \alpha + 2 (H - h) \tan \alpha \\
& = (2H -h) \tan \alpha
 \end{align*}
由于 $\tan \alpha = \dfrac{\sin \alpha}{ \sqrt{1 - \sin^{2} \alpha}}$，利用方程~\eqref{eq-193} 得
\begin{equation}%
L = \frac{2H - h}{\sqrt{n^{2}-1}}
\label{eq-194}
\end{equation}
代入 $n=4/3$ 后，得
\begin{equation*}%
L = \frac{3}{\sqrt{7}}(2H -h)
\end{equation*}
\tchr 完全正确。那么潜水员朝头顶方向看去，会看到怎样的图景呢？

\stda 在正上方他会看到一个发光的圆，其半径为
\begin{equation*}%
l = \frac{(H-h)}{\sqrt{n^{2}-1}} = \frac{3}{\sqrt{7}} (H-h)
\end{equation*}
（见\figr{fig-130}）。在这个圆的范围之外，他看到的是位于湖底的物体的像。
\pnote{这个头顶的亮圆常称“斯涅耳窗”（Snell's window）。由式~\eqref{eq-193} 得临界角 $\alpha\approx48.6^{\circ}$：潜水员透过半顶角约 $48.6^{\circ}$ 的圆锥看到水面上的整个天空，圆外则是湖底物体经全反射所成的像。人教版选择性必修第一册第四章第 2 节《全反射》讨论的正是临界角与这类现象。}

\stdb 如果潜水员所站的那部分湖底不是水平的，而是倾斜的，会怎样呢？

\tchr 这时距离 $L$ 显然将依赖于潜水员观看的方向。你不难看出：当潜水员沿倾斜的湖底朝上坡方向看时，这个距离最小；朝相反方向看时最大。上一题所得的结果，现在只有当潜水员沿湖深不变的方向（平行于湖岸）看时才能应用。倾斜湖底的问题将作为家庭作业给出（见习题 74）。

\stda 在光束的路径上插入一组两面平行的透明平板（plane-parallel plates），能改变光束的方向吗？

\tchr 你们怎么看？

\stda 我认为原则上可以。我们知道，光束在板内被折射后，是沿另一个方向传播的。

\stdb 我不同意。光束从板中射出后，仍将平行于原来的方向。

\tchr 那就请你用几块折射率不同的板组成的系统来证明这一点。

\stdb 我取三块板，折射率分别为 $n_{1}$、$n_{2}$ 和 $n_{3}$。光束通过这个系统的路径如\figr{fig-131} 所示。对光束在每个界面上的折射，可以写出
\begin{align*}
\frac{\sin \alpha_{0}}{\sin \alpha_{1}}  = n_{1}; \quad \frac{\sin \alpha_{1}}{\sin \alpha_{2}} = \frac{n_{2}}{n_{1}};
\quad  \frac{\sin \alpha_{2}}{\sin \alpha_{3}} = \frac{n_{3}}{n_{2}}; \quad \frac{\sin \alpha_{3}}{\sin \alpha_{4}}  = \frac{1}{n_{3}};
\end{align*}
\begin{figure}[!h]%
\centering
\includegraphics[width=0.4\textwidth]{figs/ch-31/fig-131a}
\caption[光束通过折射率不同的平板。]{光束通过几块折射率不同的平板时发生的折射。\label{fig-131}}
\end{figure}
把这些等式的左边与左边、右边与右边分别相乘，得
\begin{equation*}
\frac{\sin \alpha_{0}}{ \sin \alpha_{4}} =1
\end{equation*}
于是 $\alpha_{0} = \alpha_{4}$，这正是我们要证明的。

\tchr 完全正确。现在我们来讨论几何光学（geometrical optics）定律的适用范围。

\stdb 这些定律不适用于与所涉及光的波长（wavelength）同数量级以及更短的距离。在这样短的距离上，光的波动性开始显现。
\pnote{波长量级的限制即光的衍射，人教版选择性必修第一册第四章《光的衍射》一节讨论的就是它；几何光学可看作波动光学在波长趋于零时的极限。至于大距离一头的限制，则是接下来这段对话的特色内容，一般中学教材不展开。}

\tchr 你说得对。这一点应试者通常都理解得相当好。你能不能说说，几何光学定律的适用性从另一头——大距离那一头——看，会受到什么限制？

\stdb 只要距离大于光的波长，光就可以在几何光学的范围内来考虑。至少以前别人是这样对我们讲的。我认为，对大距离使用几何光学没有任何限制。

\tchr 你错了。请想象这样一幅图景：你向太空发出一束光，并且完全排除它散射的可能。设在 1 秒钟内，你把发射光束的装置转过 \ang{60}。问题是：在转动过程中，在离装置超过 300000 km 的距离上，光束的速度是多大？

\stdb 我明白您的问题。这样的点必然以大于光速的速度移动。不过，按照相对论，大于光速的速度只有当它是实物物体的速度时才不可能实现。而这里我们处理的是一束光。

\tchr 怎么，难道光束不是物质的吗？由此可见，距离过大时，几何光学便不自洽了。这里必须考虑到：光束是光的粒子流，这种粒子称为光子（photon）。在装置转动之前已经发射出去的那些光子，对随后的转动“一无所知”，仍沿着它们被发射时的方向继续传播；新的光子则沿新的方向发射。因此，我们观察不到整束光的转动。
\pnote{核算一下：转动角速度为 \ang{60}/s，即 $\omega=\pi/3~\mathrm{rad/s}$；在 $R=3\times10^{8}~\mathrm{m}$（300000 km，即 1 光秒）处，“光斑”的移动速度 $v=\omega R\approx3.1\times10^{8}~\mathrm{m/s}$，确实超过真空光速 $c$。这并不违反相对论：相对论禁止的是物质、能量或信号以超过 $c$ 的速度传播，而光斑只是先后到达不同位置的光子照出的效果——没有一颗光子沿光斑的路径移动，也不能借助光斑超光速传递信息。学生乙想用“光束不是实物”为超光速开脱，教师的光子图像给出了正确的回答：沿各自路径传播的每一颗光子本身仍以光速 $c$ 运动。}

\stdb 那么怎样定量地估计几何光学定律在大距离一头的适用界限呢？

\tchr 这些距离应当满足：光通过它们所需的时间，远小于给定问题中的任何特征时间（例如，远小于发射光束的装置转动所需的时间）。这时，光束作为整体不会被破坏，我们就可以放心地运用几何光学的定律了。


\end{dialogue}

\section*{习题}
\label{prob-31}
\addcontentsline{toc}{section}{习题}

\begin{enumerate}[resume*=problems]
\item 我们从正上方竖直地观看一个被玻璃板盖住、位于水下的物体。玻璃板厚 \SI{5}{\centi\meter}，板上覆有一层 \SI{10}{\centi\meter} 厚的水。玻璃的折射率为 1.6。我们看到物体的像距水面多远？

\item 一名身高 \SI{1.8}{\meter} 的潜水员，站在湖中水深 \SI{5}{\meter} 处的湖底上。湖底是一个倾斜 \ang{15} 角的平面。求沿湖底从他所站之点到他能经水面反射看到的湖底点的最小距离。

\item 有一块厚 \SI{5}{\centi\meter}、折射率为 1.5 的玻璃板。光从空气以多大的入射角射到板上时，被板反射的光线与折射光线互相垂直？并对这一入射角计算光线通过玻璃板所产生的侧移（lateral displacement）。

\item 有一块厚度为 $d$、折射率为 $n$ 的玻璃板。从空气射到板上的光线，其入射角等于制作该板的玻璃的全反射临界角。计算光线通过此板所产生的侧移。

\end{enumerate}

\cleardoublepage
% NEWTERMS: law of reflection / law of refraction | 反射定律 / 折射定律 | 光学首章定名，ch-32/33（未译）沿用；初中已学（人教版八上第四章），高中折射定律见选择性必修第一册 4.1
% NEWTERMS: angle of incidence / angle of reflection / angle of refraction | 入射角 / 反射角 / 折射角 | |
% NEWTERMS: normal (to the boundary) | 法线 | 入射点处所作的界面法线 |
% NEWTERMS: index of refraction | 折射率 | 原书符号 $n$；人教版定义 $n=c/v$；式 (eq-189) 右端 $n_2/n_1$ 即相对折射率
% NEWTERMS: total internal reflection | 全反射 | 人教版选择性必修第一册 4.2《全反射》
% NEWTERMS: critical angle | 临界角 | 即全反射临界角
% NEWTERMS: geometrical optics | 几何光学 | 首现括注；本章讨论其双向适用界限 |
% NEWTERMS: wavelength | 波长 | 首现括注
% NEWTERMS: photon | 光子 | 首现括注；光束＝光子流
% NEWTERMS: plane-parallel plate | 平行平面板（两面平行的板） | 首现括注；侧移问题见习题 75、76
% NEWTERMS: lateral displacement | 侧移 | 习题 75 首现括注
% NEWTERMS: Snell's window | 斯涅耳窗 | 仅评注层引入（潜水员头顶亮圆）
% NEWTERMS: light ray / (light) beam | 光线 / 光束 | 常用词，glossary 此前未收
% NEWTERMS: scattering | 散射 | 思想实验中“完全排除散射”语境
% NEWTERMS: image | 像 | 本章均为观察到的（虚）像
